\documentclass[a4paper,11pt]{article}
\usepackage{fontspec}
\usepackage{polyglossia}
\setmainlanguage{vietnamese}
\setmainfont{DejaVu Serif}
\setsansfont{DejaVu Sans}
\setmonofont{DejaVu Sans Mono}
\usepackage[margin=2.2cm]{geometry}
\usepackage{titlesec}
\usepackage{enumitem}
\usepackage{longtable}
\usepackage{booktabs}
\usepackage{array}
\usepackage{xcolor}
\usepackage{hyperref}
\usepackage{parskip}
\usepackage{seqsplit}

\hypersetup{
  colorlinks=true,
  linkcolor=blue,
  urlcolor=blue
}

\titleformat{\section}{\Large\bfseries}{\thesection.}{0.6em}{}
\titleformat{\subsection}{\large\bfseries}{\thesubsection}{0.6em}{}

\newcolumntype{L}[1]{>{\raggedright\arraybackslash}p{#1}}

\title{\textbf{Kế hoạch Giao tiếp Nhóm \& Configuration Management}\\[4pt]
\large Dự án: Hệ thống Quản lý Du thuyền (Cruise Management System)\\
\normalsize Jira: SCRUM-14 -- Project Communications \& Configuration Management}
\author{}
\date{}

\begin{document}
\maketitle
\vspace{-1.5cm}
\tableofcontents
\vspace{0.5cm}

\section{Mục tiêu}

Tài liệu này quy định cách nhóm phối hợp, trao đổi thông tin và quản lý phiên bản mã nguồn/tài liệu trong suốt vòng đời dự án, nhằm:

\begin{itemize}[leftmargin=1.2cm]
  \item Đảm bảo thông tin thông suốt giữa các thành viên phụ trách các module khác nhau.
  \item Giảm xung đột code khi nhiều người cùng làm việc trên Git.
  \item Đảm bảo có thể truy vết (traceability) giữa task trên Jira, nhánh Git và commit.
\end{itemize}

\section{Kế hoạch Giao tiếp Nhóm}

\subsection{Kênh giao tiếp}

\begin{longtable}{L{4.2cm} L{5.3cm} L{3.5cm}}
\toprule
\textbf{Kênh} & \textbf{Mục đích} & \textbf{Đối tượng} \\
\midrule
\endhead
Zalo/Slack nhóm (kênh chính) & Trao đổi nhanh, thông báo hằng ngày, hỏi đáp kỹ thuật & Toàn team \\
Jira (comment trong task) & Trao đổi liên quan trực tiếp đến 1 task/issue cụ thể, cập nhật trạng thái & Toàn team \\
Git (commit message, Pull Request comment) & Trao đổi liên quan đến thay đổi code cụ thể & Dev liên quan + reviewer \\
Email & Thông báo chính thức, báo cáo với giảng viên/khách hàng, biên bản họp & Toàn team + bên liên quan \\
Google Meet/Zoom & Họp định kỳ, demo, giải quyết vấn đề phức tạp & Toàn team \\
LaTeX (source trên Git) & Soạn thảo và lưu trữ tài liệu thiết kế, biên bản họp, tài liệu đặc tả dưới dạng mã nguồn \texttt{.tex} & Toàn team \\
\bottomrule
\end{longtable}

\subsection{Tần suất họp}

\begin{longtable}{L{3.8cm} L{3.8cm} L{6.5cm}}
\toprule
\textbf{Loại họp} & \textbf{Tần suất} & \textbf{Nội dung} \\
\midrule
\endhead
Daily standup & Hằng ngày (15 phút, đầu giờ làm việc/tối) & Hôm qua làm gì -- hôm nay làm gì -- có vướng mắc gì \\
Sprint Planning & Đầu mỗi sprint (1--2 tuần/sprint) & Chọn task từ Backlog, ước lượng, phân công \\
Sprint Review/Demo & Cuối mỗi sprint & Demo tính năng đã hoàn thành cho cả nhóm \\
Sprint Retrospective & Cuối mỗi sprint & Rút kinh nghiệm: cái gì tốt, cái gì cần cải thiện \\
Họp đột xuất & Khi có blocker nghiêm trọng & Giải quyết vấn đề gấp ảnh hưởng tiến độ \\
\bottomrule
\end{longtable}

\subsection{Quy tắc giao tiếp}

\begin{itemize}[leftmargin=1.2cm]
  \item Mỗi thành viên cập nhật trạng thái task trên Jira (To Do $\rightarrow$ In Progress $\rightarrow$ Review $\rightarrow$ Done) \textbf{ít nhất 1 lần/ngày}.
  \item Khi gặp blocker, báo ngay trên kênh chat chính và tag người liên quan, không chờ đến buổi họp.
  \item Biên bản họp (đặc biệt Sprint Planning/Retro) được viết bằng LaTeX, commit vào thư mục \texttt{docs/} trong repository để tra cứu sau này.
  \item Phản hồi tin nhắn liên quan công việc trong vòng 24h (giờ làm việc).
\end{itemize}

\section{Configuration Management}

\subsection{Quản lý version tài liệu bằng LaTeX}

\begin{itemize}[leftmargin=1.2cm]
  \item Toàn bộ tài liệu (đặc tả yêu cầu, thiết kế CSDL, biên bản họp...) được soạn thảo bằng \textbf{LaTeX} (\texttt{.tex}) và lưu trong repository Git, thay vì Google Docs/Word, để version được quản lý cùng cơ chế với code (commit, diff, branch, PR).
  \item Cấu trúc thư mục đề xuất trong repo:
\begin{verbatim}
docs/
├── dac-ta-yeu-cau/       # Dac ta yeu cau
├── thiet-ke-csdl/        # ERD, thiet ke CSDL
├── bien-ban-hop/         # Bien ban hop theo tung sprint
└── shared/               # File .sty, hinh anh, template dung chung
\end{verbatim}
  \item Mỗi tài liệu là một file \texttt{.tex} riêng; ảnh minh họa/sơ đồ đặt trong thư mục con \texttt{images/} tương ứng.
  \item Đặt tên file theo quy ước: \texttt{ten-tai-lieu\_vX.Y.tex} (ví dụ: \texttt{erd-cruise-system\_v1.2.tex}). Số phiên bản (X.Y) tăng theo mỗi lần thay đổi nội dung đáng kể; lịch sử chi tiết đã có sẵn qua \texttt{git log}/\texttt{git blame} nên không cần ghi ngày tháng trong tên file.
  \item Commit message cho tài liệu tuân theo cùng quy ước với code: \texttt{SCRUM-XX: mô tả thay đổi tài liệu} (xem mục 3.2).
  \item Build ra PDF bằng \texttt{pdflatex}/\texttt{latexmk}; file PDF \textbf{không} commit vào repo (thêm vào \texttt{.gitignore}), chỉ build khi cần chia sẻ bản đọc hoặc đính kèm vào Jira/Drive tạm thời để nộp báo cáo.
  \item Thay đổi lớn ở tài liệu quan trọng (đặc tả, ERD) phải mở Pull Request để cả nhóm review nội dung trước khi merge, và thông báo trên kênh chat chính.
\end{itemize}

\subsection{Quản lý code qua Git/Jira}

\begin{itemize}[leftmargin=1.2cm]
  \item Repository chính đặt trên \textbf{GitHub/GitLab}, liên kết với Jira để mỗi commit/PR có thể trỏ về đúng issue (dùng mã issue, ví dụ \texttt{SCRUM-21}, trong tên nhánh và commit message).
  \item \textbf{Commit message} theo dạng: \texttt{SCRUM-21: mô tả ngắn gọn thay đổi}.
  \item \textbf{Pull Request} bắt buộc có ít nhất 1 người review trước khi merge vào \texttt{develop}/\texttt{main}.
  \item Nhánh \texttt{main} luôn ở trạng thái ổn định, chỉ merge từ \texttt{develop} sau khi test.
\end{itemize}

\subsection{Quy ước đặt nhánh (branch naming) theo module}

Cấu trúc chung:
\begin{verbatim}
<loai>/<module>/<ma-issue>-<mo-ta-ngan>
\end{verbatim}

\begin{longtable}{L{3.5cm} L{9.5cm}}
\toprule
\textbf{Loại nhánh} & \textbf{Khi dùng} \\
\midrule
\endhead
\texttt{feature/...} & Phát triển tính năng mới \\
\texttt{bugfix/...} & Sửa lỗi \\
\texttt{hotfix/...} & Sửa lỗi khẩn cấp trên production \\
\texttt{refactor/...} & Tái cấu trúc code, không đổi chức năng \\
\bottomrule
\end{longtable}

\textbf{Ví dụ cụ thể cho module của Trọng Hải (Hành khách \& Đăng ký):}

\begingroup
\small
\begin{longtable}{L{3.8cm} L{9.2cm}}
\toprule
\textbf{Bảng CSDL phụ trách} & \textbf{Ví dụ tên nhánh} \\
\midrule
\endhead
Passenger & \texttt{\seqsplit{feature/passenger/SCRUM-21-quan-ly-thong-tin-hanh-khach}} \\
Booking & \texttt{\seqsplit{feature/booking/SCRUM-22-quan-ly-booking}} \\
Cabin & \texttt{\seqsplit{feature/cabin/SCRUM-23-quan-ly-cabin}} \\
ActivityRegistration & \texttt{\seqsplit{feature/activity-registration/SCRUM-24-dang-ky-hoat-dong}} \\
ExcursionRegistration & \texttt{\seqsplit{feature/excursion-registration/SCRUM-25-dang-ky-tour}} \\
CheckIn & \texttt{\seqsplit{feature/checkin/SCRUM-26-checkin-qr-rfid}} \\
\bottomrule
\end{longtable}
\endgroup

\noindent\textit{Áp dụng cùng quy ước cho các module khác của nhóm (ví dụ: \texttt{feature/finance/...}, \texttt{feature/crew/...}, \texttt{feature/schedule/...} tùy theo phân công thực tế).}

\end{document}
